%
% conversation-mode-call-state.tex
%
% Z Specification for the Conversation Mode call state machine.
% Formal model of idle/listening/waiting/speaking, turn detection, and
% barge-in -- written for Slice 1a of the Conversation Mode Phase 1 epic
% (docs/conversation-mode-prd.tex Chapter 2, section "The unresolved core:
% session attachment", and Chapter 3's testability addendum).
%

\documentclass[a4paper,10pt,fleqn]{article}
\usepackage[margin=1in]{geometry}
\usepackage{fuzz}
\usepackage[colorlinks=true,linkcolor=blue,citecolor=blue,urlcolor=blue]{hyperref}

\begin{document}

\title{Vox Conversation Mode: A Z Specification for the Call State Machine}
\author{Formal Model of idle/listening/waiting/speaking}
\date{August 2026}
\hypersetup{
  pdftitle={Vox Conversation Mode Call State Machine: A Z Specification},
  pdfauthor={Punt Labs},
  pdfsubject={Z Specification},
  pdfcreator={fuzz/probcli}
}
\maketitle

\tableofcontents
\newpage

\section{Introduction}

A Conversation Mode call moves through four modes -- \emph{idle} (no call),
\emph{listening} (waiting to hear the human), \emph{waiting} (a turn has been
sent to the agent; no reply audio yet), and \emph{speaking} (the agent's reply
is playing). This document models that machine formally, per
\texttt{docs/WORKFLOW.md}'s stateful-audio gate: the machine has four modes
and transitions between them, invariants must hold across those transitions
(a pending human utterance is never silently dropped), and a wrong transition
here is exactly the class of defect that gate exists to catch before code is
written, not after.

Two questions were named in \texttt{docs/conversation-mode-prd.tex} as
requiring resolution \emph{before} this model was written, not discovered
while writing it (Chapter 2, ``Known unknowns''; Chapter 3,
\S\ref{sec:test-di} item 6). Both are resolved here, explicitly:

\begin{enumerate}
  \item \textbf{Which detector governs a human speaking during the agent's
    silent think-time (\emph{waiting}, before any reply audio starts)?}
    Resolved in \S\ref{sec:detector} as a totally-defined function from
    \texttt{Mode} to \texttt{Detector}: the \emph{turn detector} governs
    both \emph{listening} and \emph{waiting} -- barge-in detection requires
    agent audio to interrupt, and no agent audio exists until \emph{speaking}
    begins, so there is nothing for a barge-in detector to do during
    \emph{waiting}. A human who speaks again during \emph{waiting} is not
    interrupting anything; the turn detector's normal end-of-utterance logic
    applies, and what it captures is held as a \emph{pending addendum}
    (\S\ref{sec:addendum}), not discarded and not immediately dispatched as a
    second, concurrent turn (FR-4: one call is one human and one agent
    session at a time; the current turn is still in flight).
  \item \textbf{Where does the single serialized dispatch point for this
    state machine live?} Resolved in \S\ref{sec:dispatch} as a design
    requirement on the machine's \emph{caller}, not a property the pure
    state schema below can enforce by itself -- exactly as
    \texttt{voxd/programs/program.py}'s \texttt{Program} already documents
    for its own single-writer assumption. \S\ref{sec:dispatch} states the
    resolution and why it is the right one for this codebase.
\end{enumerate}

This model does not attempt session attachment, turn-detection acoustics, or
synthesis timing -- those are described informally in
\texttt{docs/conversation-mode-prd.tex} Chapter 2 and are the subject of
Slice 1b onward. This model's job is the state machine those components
drive: the shape \texttt{Program} already proved out for Audio Programs
(\texttt{docs/audio-programs.tex}), applied to a live call instead of a
playlist.

\section{Basic Types}

There is no session or agent identifier in this model. FR-4 requires
exactly one human and one Claude Code agent session per call, and that
constraint is structural (one \texttt{CallState} instance per call, created
by \texttt{StartCall} and destroyed by the return to \texttt{idle}) rather
than a value carried in the state, the same reasoning
\texttt{docs/audio-programs.tex} gives for \texttt{Program} carrying no
owner.

\begin{zed}
  [TURN]
\end{zed}

$TURN$ is an opaque handle for one transcribed human utterance -- what a
\texttt{TranscribedTurn} carries in the implementation
(\texttt{src/punt\_vox/voxd/conversation\_mode/turn.py}). Its
internal structure (the transcript text, a confidence signal) is outside
this model; the state machine only needs to know a turn was detected, not
what it says.

\section{Free Types}

\begin{zed}
  Mode ::= idle | listening | waiting | speaking
\end{zed}

\begin{zed}
  Detector ::= noDetector | turnDetector | bargeInDetector
\end{zed}

\begin{zed}
  ZBOOL ::= ztrue | zfalse
\end{zed}

A Z boolean (avoids the ProB/B keyword clash with \texttt{BOOL}), following
the convention already established in \texttt{docs/audio-programs.tex}.

\section{Which Detector Governs Each Mode}
\label{sec:detector}

\begin{axdef}
  activeDetector : Mode \fun Detector
  \where
  activeDetector~idle = noDetector \\
  activeDetector~listening = turnDetector \\
  activeDetector~waiting = turnDetector \\
  activeDetector~speaking = bargeInDetector
\end{axdef}

This is the resolution to gap (a): a \emph{total} function, so there is no
mode for which ``which detector'' is left unanswered. \texttt{waiting} maps
to \texttt{turnDetector}, the same detector as \texttt{listening} -- the two
non-speaking modes share one detector, and only \texttt{speaking} switches to
\texttt{bargeInDetector}. This is not an arbitrary tie-break: barge-in is
defined (FR-8) as interrupting the agent's \emph{speech}, and there is no
speech to interrupt before \texttt{speaking} begins. Modelling
\texttt{activeDetector} as a total function, rather than leaving the
\texttt{waiting} case as prose, is what makes the resolution checkable --
\texttt{fuzz} rejects the specification if the function is not total over
\texttt{Mode}.

\section{Pending Addenda}
\label{sec:addendum}

\begin{schema}{CallState}
  mode : Mode \\
  hasPendingAddendum : ZBOOL
  \where
  hasPendingAddendum = ztrue \implies mode \in \{waiting, speaking\}
\end{schema}

The invariant is the second half of gap (a)'s resolution: a pending addendum
--- speech the turn detector captured during \texttt{waiting} that was not
yet part of any turn sent to the agent --- can only exist while the call has
not yet returned to \texttt{listening}. Read contrapositively: whenever the
call is in \texttt{idle} or \texttt{listening}, \texttt{hasPendingAddendum}
is \texttt{zfalse} by construction. Every operation that lands on
\texttt{listening} (\texttt{BargeIn}, \texttt{ReplyEnds}) must therefore
discharge the flag explicitly, in the same operation -- the model will not
type-check otherwise, because the invariant would be violated on entry to
\texttt{listening}. This is a formal way of saying ``an addendum is never
silently dropped'': it cannot survive a transition back to
\texttt{listening} without the operation that makes that transition
explicitly accounting for it. What the implementation does with the
discharged content (fold it into the start of the next turn, per the same
principle FR-9 already establishes for barge-in speech) is outside this
model's scope, but the model pins the \emph{moment} that hand-off must
happen: exactly at \texttt{BargeIn} or \texttt{ReplyEnds}, never later.

\section{Initialisation}

\begin{schema}{Init}
  CallState'
  \where
  mode' = idle \\
  hasPendingAddendum' = zfalse
\end{schema}

\section{Operations}

\subsection{Call lifecycle}

\begin{schema}{StartCall}
  \Delta CallState
  \where
  mode = idle \\
  mode' = listening \\
  hasPendingAddendum' = zfalse
\end{schema}

A call begins listening immediately (UC-1, FR-1): there is no
intermediate ``setting up'' mode between \texttt{idle} and
\texttt{listening}.

\begin{schema}{EndCall}
  \Delta CallState
  \where
  mode \neq idle \\
  mode' = idle \\
  hasPendingAddendum' = zfalse
\end{schema}

An explicit hangup (UC-9, FR-2) ends the call from any active mode,
including mid-\texttt{speaking} -- the human can end the call while the
agent is still talking.

\begin{schema}{TimeoutCall}
  \Delta CallState
  \where
  mode = listening \\
  mode' = idle \\
  hasPendingAddendum' = zfalse
\end{schema}

The bounded-inactivity timeout (FR-2) fires only from \texttt{listening}:
that is the one mode in which the human is expected to be about to speak
and nothing is otherwise in flight. \texttt{waiting} and \texttt{speaking}
have their own bounded-time failure paths (a synthesis call that never
returns, Chapter 2's ``Known unknowns'') that are a distinct concern from
mic-inactivity and are out of this model's scope, per the PRD's own framing
of that gap as unresolved.

\subsection{Turn detection}

\begin{schema}{TurnDetected}
  \Delta CallState
  \where
  mode = listening \\
  mode' = waiting \\
  hasPendingAddendum' = zfalse
\end{schema}

The turn detector (\S\ref{sec:detector}: active in \texttt{listening}) fires
end-of-turn (FR-5); the transcribed turn is sent to the agent session
(\texttt{SessionAttach.send\_turn}, out of this model's scope) and the call
enters \texttt{waiting} -- no reply audio yet, but a turn is in flight.

\begin{schema}{CaptureDuringWait}
  \Delta CallState
  \where
  mode = waiting \\
  mode' = waiting \\
  hasPendingAddendum' = ztrue
\end{schema}

This is the operation gap (a) exists to name: the turn detector
(\S\ref{sec:detector}: still active in \texttt{waiting}) fires again while
the call is already waiting on the agent's reply. FR-4 rules out
dispatching a second, concurrent turn to the same session -- one call is one
human and one agent session, and the first turn has not yet been answered --
so the newly-detected speech is recorded as pending, per
\S\ref{sec:addendum}, rather than sent immediately or discarded.

\subsection{The agent's turn}

\begin{schema}{ReplyBegins}
  \Delta CallState
  \where
  mode = waiting \\
  mode' = speaking \\
  hasPendingAddendum' = hasPendingAddendum
\end{schema}

The reply's first complete portion is ready to play (FR-11); the call
enters \texttt{speaking}. A pending addendum, if any, carries through
unchanged -- it is not yet discharged, because the call has not yet
returned to \texttt{listening} (\S\ref{sec:addendum}).

\begin{schema}{BargeIn}
  \Delta CallState
  \where
  mode = speaking \\
  mode' = listening \\
  hasPendingAddendum' = zfalse
\end{schema}

The barge-in detector (\S\ref{sec:detector}: active in \texttt{speaking})
fires; the agent's audio stops (FR-8) and the call returns to
\texttt{listening}. Any pending addendum accumulated during the prior
\texttt{waiting} phase, plus the interrupting speech itself, are both
discharged here (\S\ref{sec:addendum}) and become the start of the next turn
(FR-9) -- the model does not distinguish ``lost because of the barge-in''
from ``carried forward,'' because the invariant forbids the former: reaching
\texttt{listening} with \texttt{hasPendingAddendum} still \texttt{ztrue} is
not a reachable state.

\begin{schema}{ReplyEnds}
  \Delta CallState
  \where
  mode = speaking \\
  mode' = listening \\
  hasPendingAddendum' = zfalse
\end{schema}

The reply finishes without interruption; the call signals ``ready to
listen'' audibly (FR-3, NFR-6) and returns to \texttt{listening}. Any
pending addendum is discharged the same way as in \texttt{BargeIn} -- the
two operations are the only two paths into \texttt{listening} from an active
call, and both carry the same obligation.

\subsection{Observation}

\begin{schema}{CurrentDetector}
  \Xi CallState \\
  detector! : Detector
  \where
  detector! = activeDetector~mode
\end{schema}

A non-mutating query: which detector is live right now. This is the
queryable form of gap (a)'s resolution -- a test or a caller can ask
\texttt{CurrentDetector} in any reachable state and get a definite answer,
never ``undefined for this mode.''

\section{The Single Serialized Dispatch Point}
\label{sec:dispatch}

This is the resolution to gap (b), and it is stated in prose rather than as
a further Z schema, for the same reason
\texttt{voxd/programs/program.py}'s single-writer assumption is a docstring
note on \texttt{Program} rather than a schema in
\texttt{docs/audio-programs.tex}: Z models the shape of state and the
preconditions/postconditions of a transition, not a scheduling discipline
between concurrent callers. The state schema above is safe to reason about
in isolation \emph{only} if exactly one caller applies one operation at a
time -- the same caveat \texttt{Program} carries and does not enforce
internally.

\texttt{docs/conversation-mode-prd.tex} names three components that can call
into this state machine concurrently: turn detection, barge-in detection,
and sentence-streamed synthesis (which drives \texttt{ReplyBegins}). All
three are independently-running, continuous processes for the life of a
call (Chapter 2, ``Component overview''). Copying \texttt{Program}'s pure
class shape without also copying its single-writer discipline reproduces a
data race, not a design -- this section is where that discipline is pinned
down for Conversation Mode specifically.

\textbf{Resolution: an actor with an internal, single-consumer command
queue}, not an explicit lock. \texttt{voxd} is fully async
(\texttt{docs/conversation-mode-prd.tex}, ``Sentence-streamed synthesis: one
pipeline, not two'' -- every handler in \texttt{voxd/*.py} is \texttt{async
def}), and Python's asyncio event loop is already single-threaded and
cooperative: a lock is the wrong tool for serializing access between
coroutines running on one loop, because nothing preempts a coroutine
mid-statement the way a thread can be preempted. An explicit lock would add
ceremony (acquire/release, and the risk of holding it across an
\texttt{await} and reintroducing exactly the race it exists to prevent) for
no safety a queue does not already give for free.

The concrete shape: a \texttt{CallActor} owns one \texttt{CallState} and one
\texttt{asyncio.Queue} of commands (one variant per operation above -- a
tagged union, not a bare callable, so the actor's dispatch loop can pattern
match and each command remains a first-class, testable value per PY-DP-3).
Turn detection, barge-in detection, and the synthesis pipeline each hold a
reference to the queue, never to the \texttt{CallState} or the
\texttt{Program}-shaped class directly -- they enqueue a command and return;
they never call an operation method on the state class themselves. Exactly
one task drains the queue, one command at a time, applying the matching
operation and awaiting nothing while doing so (the operation methods
themselves, mirroring \texttt{Program}, are synchronous bare-attribute
assignments with no \texttt{await} inside them) -- so no other command can
interleave mid-operation. This is the queue analogue of
\texttt{Program}'s documented (but unenforced) single-writer assumption,
made structurally impossible to violate rather than merely stated: there is
exactly one reader of the queue, and the state class itself is never handed
to any of the three concurrent producers.

This is a design requirement on Slice 1b's implementation, not a component
this design mission builds -- consistent with this mission's boundary
(the fake and Protocol under
\texttt{src/punt\_vox/voxd/conversation\_mode/} model \emph{session
attachment}, not the call actor). It is recorded here, in the Z
specification, because \texttt{docs/conversation-mode-prd.tex} explicitly
requires this document to be the place that gap is settled, not an
implementation detail discovered afterward.

\section{Known Limitations}

This model does not address two items \texttt{docs/conversation-mode-prd.tex}
names as separately unresolved, and neither was in this mission's scope to
resolve:

\begin{itemize}
  \item \textbf{A stuck or free-running detector.} The model assumes
    \texttt{activeDetector} always eventually fires the transition it
    governs. A detector that hangs (no transition ever fires) or free-runs
    (fires spuriously with no genuine speech) is a runtime failure mode
    this state machine cannot detect from the inside -- consistent with
    \texttt{docs/conversation-mode-prd.tex}'s own framing of this as an open
    question, not a gap this document was asked to close.
  \item \textbf{A synthesis call that never returns mid-turn.} The model
    assumes \texttt{ReplyBegins} is always eventually reachable from
    \texttt{waiting}. A hung synthesis call is, like the stuck detector
    above, a bounded-time failure this document does not attempt to model;
    \texttt{docs/conversation-mode-prd.tex} suggests a timeout ending the
    turn (mirroring FR-18) as the likely mechanism, but that is
    implementation, not this model's job.
\end{itemize}

\end{document}
